\begin{abstract}
Zero-shot learning (ZSL) nhằm mục đích nhận dạng các lớp chưa từng thấy mà không có các mẫu huấn luyện được gán nhãn bằng cách tận dụng các bộ mô tả ngữ nghĩa cấp lớp như các thuộc tính. Một thách thức cơ bản trong ZSL là sự sai lệch ngữ nghĩa (semantic misalignment), trong đó thông tin không liên quan đến ngữ nghĩa có trong các đặc trưng thị giác gây ra sự mơ hồ cho tương tác thị giác-ngữ nghĩa. Không giống như các phương pháp hiện có giúp triệt tiêu thông tin không liên quan đến ngữ nghĩa sau xử lý (post hoc) trong không gian đặc trưng hoặc không gian mô hình, chúng tôi đề xuất giải quyết vấn đề này ở giai đoạn đầu vào, ngăn chặn các bản vá không liên quan đến ngữ nghĩa lan truyền qua mạng. Để đạt được mục tiêu này, chúng tôi giới thiệu Semantically Contextualized Visual Patches (SVIP) cho ZSL, một framework dựa trên transformer được thiết kế để tăng cường sự liên kết thị giác-ngữ nghĩa. Cụ thể, chúng tôi đề xuất một cơ chế lựa chọn bản vá tự giám sát (self-supervised) để học cách xác định các bản vá không liên quan đến ngữ nghĩa trong không gian đầu vào một cách ưu tiên. Cơ chế này được huấn luyện với sự giám sát từ các điểm chú ý được tổng hợp trên tất cả các lớp transformer, ước tính điểm ngữ nghĩa của mỗi bản vá. Vì việc loại bỏ các bản vá không liên quan đến ngữ nghĩa khỏi chuỗi đầu vào có thể làm phá vỡ cấu trúc đối tượng, chúng tôi thay thế chúng bằng các nhúng bản vá có thể học được (learnable patch embeddings). Với việc khởi tạo từ các nhúng từ (word embeddings), chúng tôi có thể đảm bảo chúng vẫn có ý nghĩa ngữ nghĩa trong suốt quá trình trích xuất đặc trưng. Các thí nghiệm sâu rộng trên các bộ dữ liệu ZSL tiêu chuẩn đã chứng minh rằng SVIP đạt được kết quả hiệu suất hiện đại (state-of-the-art) đồng thời cung cấp các biểu diễn đặc trưng có khả năng diễn giải tốt hơn và giàu ngữ nghĩa hơn. Mã nguồn có sẵn tại \url{https://github.com/uqzhichen/SVIP}.
\end{abstract}
